\chapter{附录A：代数基础 习题}
\difficultykey

\section{基础题 \easy}

\begin{problem}{线性无关组的延拓}
\easy
设 $V$ 是有限维向量空间，$S=\{v_1,\dots,v_r\}$ 是一个线性无关组。证明 $S$ 可以延拓为 $V$ 的一组基。
\end{problem}
\begin{solution}
若 $S$ 已经张成 $V$，它本身就是一组基。

若 $\langle S\rangle\neq V$，就取 $v_{r+1}\in V\setminus\langle S\rangle$。此时 $S\cup\{v_{r+1}\}$ 仍线性无关；否则若有
\[
a_1v_1+\cdots+a_rv_r+a_{r+1}v_{r+1}=0
\]
且 $a_{r+1}\neq0$，则 $v_{r+1}\in\langle S\rangle$，矛盾。

不断重复这一过程。因为 $V$ 有限维，至多加上有限个向量便会张成全空间，于是得到一组基。
\end{solution}

\begin{problem}{核与像的计算}
\easy
设线性映射
\[
T:\R^3\to\R^2,
\qquad
T(x,y,z)=(x+y,\,y+z).
\]
求 $\ker T$ 与 $\operatorname{im}T$，并验证秩亏定理。
\end{problem}
\begin{solution}
解方程
\[
T(x,y,z)=(0,0)
\]
即
\[
x+y=0,
\qquad y+z=0.
\]
故
\[
x=-y,
\qquad z=-y,
\]
于是
\[
\ker T=\{(-t,t,-t):t\in\R\}=\langle(-1,1,-1)\rangle.
\]
所以 $\dim\ker T=1$。

对任意 $(u,v)\in\R^2$，取 $y=0,x=u,z=v$，则 $T(u,0,v)=(u,v)$，故
\[
\operatorname{im}T=\R^2,
\qquad \dim\operatorname{im}T=2.
\]
因此
\[
\dim\ker T+\dim\operatorname{im}T=1+2=3=\dim\R^3,
\]
秩亏定理成立。
\end{solution}

\begin{problem}{特征值与特征向量}
\easy
设
\[
A=\begin{pmatrix}2&1\\1&2\end{pmatrix}.
\]
求 $A$ 的特征值与对应特征向量。
\end{problem}
\begin{solution}
特征多项式为
\[
\det(A-\lambda I)=
\begin{vmatrix}2-\lambda&1\\1&2-\lambda\end{vmatrix}
=(2-\lambda)^2-1=\lambda^2-4\lambda+3.
\]
故特征值为
\[
\lambda=1,\quad 3.
\]

当 $\lambda=3$ 时，解
\[
(A-3I)v=0
\]
得 $x=y$，可取特征向量 $(1,1)^\top$。

当 $\lambda=1$ 时，解
\[
(A-I)v=0
\]
得 $x=-y$，可取特征向量 $(1,-1)^\top$。
\end{solution}

\begin{problem}{不同特征值的对角化}
\easy
证明：若线性变换 $T:V\to V$ 在 $n$ 维向量空间上有 $n$ 个互不相同的特征值，则 $T$ 可对角化。
\end{problem}
\begin{solution}
先证明属于不同特征值的特征向量线性无关。设
\[
\lambda_1,\dots,\lambda_m
\]
两两不同，而
\[
a_1v_1+\cdots+a_mv_m=0
\]
其中 $Tv_i=\lambda_iv_i$。对上式施加 $T-\lambda_mI$，可把最后一项消去，得到关于 $v_1,\dots,v_{m-1}$ 的关系。归纳可知所有 $a_i=0$。

现在若 $T$ 有 $n$ 个互异特征值，就得到 $n$ 个线性无关特征向量；在 $n$ 维空间中，这些向量构成一组基。相对于这组基，$T$ 的矩阵就是对角矩阵，所以 $T$ 可对角化。
\end{solution}

\begin{problem}{子群判别}
\easy
设
\[
H=\{e,(12)(34),(13)(24),(14)(23)\}\subset S_4.
\]
证明 $H$ 是 $S_4$ 的子群，并说明它与哪个熟悉的群同构。
\end{problem}
\begin{solution}
首先 $e\in H$。其次，三个非平凡元素都满足平方为 $e$，所以逆元仍在 $H$ 中。

再看封闭性，例如
\[
(12)(34)\,(13)(24)=(14)(23),
\]
其余乘积同理仍落在 $H$ 中。因此 $H$ 是子群。

它有四个元素，三个非平凡元素都阶为 $2$，故同构于 Klein 四元群
\[
V_4\cong \Z/2\Z\times\Z/2\Z.
\]
\end{solution}

\begin{problem}{核像与第一同构定理}
\easy
设群同态
\[
\varphi:\Z\to\Z/12\Z,
\qquad
\varphi(n)=\overline n.
\]
求 $\ker\varphi$ 与 $\operatorname{im}\varphi$，并写出第一同构定理给出的同构。
\end{problem}
\begin{solution}
核由满足 $\overline n=\overline0$ 的整数构成，所以
\[
\ker\varphi=12\Z.
\]
像显然是整个 $\Z/12\Z$，即
\[
\operatorname{im}\varphi=\Z/12\Z.
\]

第一同构定理给出
\[
\Z/\ker\varphi\cong\operatorname{im}\varphi,
\]
因此
\[
\Z/12\Z\cong\Z/12\Z.
\]
这里的同构正是把整数的陪集 $n+12\Z$ 送到剩余类 $\overline n$。
\end{solution}

\begin{problem}{整数环中的理想}
\easy
证明 $\Z$ 的每个理想都形如 $n\Z$，其中 $n\ge0$ 唯一确定。
\end{problem}
\begin{solution}
设 $I\subset\Z$ 是理想。若 $I=\{0\}$，则 $I=0\Z$。

若 $I\neq\{0\}$，由良序性可取其中最小的正整数 $n$。因为 $I$ 是理想，故 $n\Z\subset I$。反过来，对任意 $a\in I$，用带余除法写成
\[
a=qn+r,
\qquad 0\le r<n.
\]
由于 $a,qn\in I$，故 $r=a-qn\in I$。由 $n$ 的最小性可知 $r=0$，所以 $a\in n\Z$。因此 $I=n\Z$。

唯一性来自：若 $n\Z=m\Z$ 且 $n,m\ge0$，则 $n\in m\Z$ 且 $m\in n\Z$，于是 $n=m$。
\end{solution}

\begin{problem}{有限域的基本性质}
\easy
设 $p$ 是素数。说明有限域 $\F_p$ 恰有 $p$ 个元素，且每个非零元素都可逆。
\end{problem}
\begin{solution}
按定义
\[
\F_p=\Z/p\Z,
\]
它的元素就是模 $p$ 的剩余类
\[
\overline0,\overline1,\dots,\overline{p-1},
\]
故共有 $p$ 个元素。

若 $\overline a\neq\overline0$，则 $p\nmid a$。因为 $p$ 是素数，所以
\[
\gcd(a,p)=1.
\]
由 Bézout 等式，存在整数 $u,v$ 使
\[
ua+vp=1.
\]
模 $p$ 化后得
\[
\overline u\,\overline a=\overline1,
\]
所以每个非零元素都有逆元。
\end{solution}

\begin{problem}{自由模与非自由模}
\easy
把 $\Z$ 看成自身上的模。说明 $\Z$ 是自由 $\Z$-模，而 $\Z/ n\Z$ 在 $n>1$ 时不是自由 $\Z$-模。
\end{problem}
\begin{solution}
$\Z$ 由元素 $1$ 生成，而且任意元素唯一写成 $m\cdot1$，所以
\[
\Z\cong \Z^1
\]
是自由秩 $1$ 的 $\Z$-模。

若 $\Z/n\Z$ 是自由 $\Z$-模，则它作为阿贝尔群必与某个 $\Z^r$ 同构。可是 $\Z^r$ 在 $r\ge1$ 时有无穷多个元素，而 $\Z/n\Z$ 只有有限多个元素；若 $r=0$ 则它是零模。故当 $n>1$ 时，$\Z/n\Z$ 不是自由 $\Z$-模。
\end{solution}

\section{进阶题 \medium}

\begin{problem}{Jordan标准形的识别}
\medium
设
\[
A=\begin{pmatrix}
2&1&0\\0&2&1\\0&0&2
\end{pmatrix}.
\]
求 $A$ 的特征多项式、最小多项式，并写出其 Jordan 标准形。
\end{problem}
\begin{solution}
$A$ 是上三角矩阵，所以特征值全是对角元 $2$。故特征多项式为
\[
\chi_A(x)=(x-2)^3.
\]

又有
\[
A-2I=
\begin{pmatrix}
0&1&0\\0&0&1\\0&0&0
\end{pmatrix},
\qquad
(A-2I)^2\neq0,
\qquad
(A-2I)^3=0.
\]
因此最小多项式为
\[
\mu_A(x)=(x-2)^3.
\]
最小多项式次数为 $3$，说明只有一个大小为 $3$ 的 Jordan 块，所以 Jordan 标准形就是它本身：
\[
J_3(2)=
\begin{pmatrix}
2&1&0\\0&2&1\\0&0&2
\end{pmatrix}.
\]
\end{solution}

\begin{problem}{Cayley和Hamilton定理的应用}
\medium
设
\[
A=\begin{pmatrix}1&2\\1&4\end{pmatrix}.
\]
利用 Cayley--Hamilton 定理求一个只含 $I$ 与 $A$ 的公式来表示 $A^{-1}$。
\end{problem}
\begin{solution}
先求特征多项式：
\[
\chi_A(x)=x^2-5x+2.
\]
由 Cayley--Hamilton 定理，
\[
A^2-5A+2I=0.
\]
把它改写为
\[
A(A-5I)=-2I.
\]
因为右边可逆，所以 $A$ 可逆，且
\[
A^{-1}=\frac{1}{2}(5I-A).
\]
直接展开可得
\[
A^{-1}=\frac12
\begin{pmatrix}
4&-2\\-1&1
\end{pmatrix},
\]
与通常求逆公式一致。
\end{solution}

\begin{problem}{Sylow定理的应用}
\medium
证明每个 $15$ 阶群都是循环群。
\end{problem}
\begin{solution}
设 $G$ 是 $15=3\cdot5$ 阶群。由 Sylow 定理，Sylow $5$-子群个数 $n_5$ 满足
\[
n_5\equiv1\pmod5,
\qquad n_5\mid3.
\]
故 $n_5=1$。同理，Sylow $3$-子群个数 $n_3$ 满足
\[
n_3\equiv1\pmod3,
\qquad n_3\mid5,
\]
故 $n_3=1$。

因此 Sylow $5$-子群 $P$ 与 Sylow $3$-子群 $Q$ 都正规，且
\[
G\cong P\times Q.
\]
又 $P\cong\Z/5\Z$，$Q\cong\Z/3\Z$ 都是循环群，并且 $3,5$ 互素，所以
\[
P\times Q\cong\Z/15\Z.
\]
故 $G$ 是循环群。
\end{solution}

\begin{problem}{中国剩余定理的具体例子}
\medium
证明
\[
\Z/12\Z\cong \Z/3\Z\times\Z/4\Z,
\]
并写出同构与逆同构。
\end{problem}
\begin{solution}
定义同态
\[
\Phi:\Z/12\Z\to\Z/3\Z\times\Z/4\Z,
\qquad
\overline x\mapsto(\overline x\bmod3,\overline x\bmod4).
\]
它是良定义的环同态。

若 $\Phi(\overline x)=(\overline0,\overline0)$，则 $3\mid x$ 且 $4\mid x$，故 $12\mid x$，所以核为零。两边都有 $12$ 个元素，因此 $\Phi$ 是同构。

逆同构可由解同余组给出。例如
\[
(a,b)\in\Z/3\Z\times\Z/4\Z
\]
送到唯一满足
\[
x\equiv a\pmod3,
\qquad x\equiv b\pmod4
\]
的 $\overline x\in\Z/12\Z$。具体地，可写成
\[
\Phi^{-1}(a,b)=4a+9b\pmod{12}.
\]
\end{solution}

\begin{problem}{商环成为域的判别}
\medium
设 $F$ 是域，$f(x)\in F[x]$ 为非常数多项式。证明商环
\[
F[x]/(f)
\]
是域当且仅当 $f$ 在 $F[x]$ 中不可约。
\end{problem}
\begin{solution}
因为 $F[x]$ 是主理想整环，所以理想 $(f)$ 是极大理想当且仅当 $f$ 不可约。

而一般交换环论告诉我们：商环 $R/I$ 是域当且仅当 $I$ 是极大理想。因此
\[
F[x]/(f)\text{ 是域}
\iff (f)\text{ 极大}
\iff f\text{ 不可约}.
\]
\end{solution}

\begin{problem}{三次多项式的分裂域}
\medium
求多项式
\[
x^3-2
\]
在 $\Q$ 上的分裂域次数，并确定其 Galois 群。
\end{problem}
\begin{solution}
由有理根定理，$x^3-2$ 在 $\Q$ 上无有理根，所以不可约。设实根为
\[
\alpha=\sqrt[3]{2}.
\]
它的三个根为
\[
\alpha,\quad \omega\alpha,\quad \omega^2\alpha,
\]
其中 $\omega=e^{2\pi i/3}$。

因此分裂域为
\[
K=\Q(\alpha,\omega).
\]
先有 $[\Q(\alpha):\Q]=3$，而 $\omega\notin\Q(\alpha)$，因为后者是实域，所以再乘一个 $2$，得到
\[
[K:\Q]=6.
\]
Galois 群阶为 $6$，它作用在三个根上，既有 $3$-循环 $\alpha\mapsto\omega\alpha$，又有复共轭给出的换位，因此
\[
\Gal(K/\Q)\cong S_3.
\]
\end{solution}

\begin{problem}{PID中自由模的子模}
\medium
证明：若 $R$ 是 PID，$M\cong R^n$ 是有限生成自由 $R$-模，则 $M$ 的每个子模都是自由模。
\end{problem}
\begin{solution}
对 $n$ 做归纳。$n=1$ 时，$M\cong R$，其子模就是理想，而 PID 中每个理想都主，因此同构于 $R$ 或 $0$，都自由。

设结论对 $n-1$ 成立。取子模 $N\subset R^n$。考虑投影到第一坐标
\[
\pi:N\to R.
\]
其像是 $R$ 的一个理想，所以形如 $dR$。若像为零，则 $N\subset0\times R^{n-1}$，由归纳假设自由。若像非零，取 $x\in N$ 使 $\pi(x)=d$ 生成该像。可证明
\[
N=Rx\oplus (N\cap (0\times R^{n-1})).
\]
第二项由归纳假设自由，因此 $N$ 自由。
\end{solution}

\begin{problem}{张量积的计算}
\medium
证明
\[
\Z/m\Z\otimes_\Z \Z/n\Z\cong \Z/\gcd(m,n)\Z.
\]
\end{problem}
\begin{solution}
先用表示
\[
\Z/m\Z\cong \Z/m\Z
\]
作为 $\Z$-模，并对短正合列
\[
\Z\xrightarrow{\times m}\Z\to\Z/m\Z\to0
\]
张量上 $\Z/n\Z$，得到
\[
\Z/n\Z\xrightarrow{\times m}\Z/n\Z\to \Z/m\Z\otimes \Z/n\Z\to0.
\]
因此张量积同构于映射 $\times m$ 的余核，即
\[
(\Z/n\Z)/m(\Z/n\Z).
\]

在循环群 $\Z/n\Z$ 中，子群 $m(\Z/n\Z)$ 的大小为
\[
\frac{n}{\gcd(m,n)},
\]
所以余核大小是 $\gcd(m,n)$。它本身又是循环群，因此
\[
\Z/m\Z\otimes_\Z\Z/n\Z\cong\Z/\gcd(m,n)\Z.
\]
\end{solution}

\begin{problem}{十二阶阿贝尔群的分类}
\medium
分类所有 $12$ 阶阿贝尔群。
\end{problem}
\begin{solution}
因为
\[
12=4\cdot3=2^2\cdot3,
\]
有限阿贝尔群分解为其 Sylow 部分直积。$3$-部分只能是
\[
\Z/3\Z,
\]
而 $2$-部分有两种可能：
\[
\Z/4\Z,
\qquad
\Z/2\Z\times\Z/2\Z.
\]
因此所有 $12$ 阶阿贝尔群恰有两类：
\[
\Z/4\Z\times\Z/3\Z\cong\Z/12\Z,
\]
和
\[
(\Z/2\Z\times\Z/2\Z)\times\Z/3\Z
\cong\Z/6\Z\times\Z/2\Z.
\]
所以共有两种，而不是更多。
\end{solution}

\section{挑战题 \hard}

\begin{problem}{特殊线性群的阶}
\hard
设 $p$ 为素数。计算
\[
|\SL_2(\F_p)|.
\]
\end{problem}
\begin{solution}
先算 $\GL_2(\F_p)$ 的阶。第一列可以是任意非零向量，共
\[
p^2-1
\]
种；第二列不能落在第一列张成的一维子空间里，因此有
\[
p^2-p
\]
种。故
\[
|\GL_2(\F_p)|=(p^2-1)(p^2-p).
\]

行列式映射
\[
\det:\GL_2(\F_p)\to\F_p^\times
\]
是满射，且像大小为 $p-1$。各纤维等大，因此
\[
|\SL_2(\F_p)|=\frac{|\GL_2(\F_p)|}{p-1}
=\frac{(p^2-1)(p^2-p)}{p-1}=p(p^2-1).
\]
\end{solution}

\begin{problem}{p群的中心与上中心列}
\hard
设 $G$ 是有限 $p$ 群。证明 $Z(G)\neq\{e\}$，并说明如何由此构造上中心列。
\end{problem}
\begin{solution}
考虑共轭作用的类方程：
\[
|G|=|Z(G)|+\sum [G:C_G(x_i)],
\]
其中和式遍历非中心共轭类代表元。每个指数
\[
[G:C_G(x_i)]
\]
都是 $p$ 的正次幂，因此被 $p$ 整除。又 $|G|$ 本身是 $p$ 的幂，所以 $|Z(G)|$ 也必须被 $p$ 整除，从而
\[
|Z(G)|\ge p.
\]
因此中心非平凡。

于是可以定义
\[
Z_0(G)=\{e\},
\qquad
Z_{i+1}(G)/Z_i(G)=Z(G/Z_i(G)).
\]
因为每一步商群仍是有限 $p$ 群，其中心非平凡，所以若 $Z_i(G)\neq G$，就有真包含
\[
Z_i(G)\subsetneq Z_{i+1}(G).
\]
有限性保证该过程在有限步后达到 $G$，这就是上中心列。
\end{solution}

\begin{problem}{双二次扩张的Galois对应}
\hard
设
\[
K=\Q(\sqrt2,\sqrt3).
\]
求 $\Gal(K/\Q)$，并写出所有中间域与子群的对应关系。
\end{problem}
\begin{solution}
扩张 $K/\Q$ 是由两个独立平方根生成的双二次扩张，因此
\[
[K:\Q]=4.
\]
每个自同构独立地决定
\[
\sqrt2\mapsto\pm\sqrt2,
\qquad
\sqrt3\mapsto\pm\sqrt3,
\]
故
\[
\Gal(K/\Q)\cong V_4.
\]

它有三个指数 $2$ 的子群，对应三个二次子域：
\begin{itemize}
\item[(i)] 固定 $\sqrt2$ 的子群，对应子域 $\Q(\sqrt2)$；
\item[(ii)] 固定 $\sqrt3$ 的子群，对应子域 $\Q(\sqrt3)$；
\item[(iii)] 固定 $\sqrt6$ 的子群，对应子域 $\Q(\sqrt6)$。
\end{itemize}
此外，平凡子群对应全域 $K$，整个 Galois 群对应底域 $\Q$。这给出完整的 Galois 对应。
\end{solution}

\begin{problem}{四次分裂域的Galois群}
\hard
求多项式
\[
x^4-2
\]
在 $\Q$ 上的分裂域、扩张次数与 Galois 群，并写出其三个二次子域。
\end{problem}
\begin{solution}
设
\[
\alpha=\sqrt[4]{2}.
\]
则四个根为
\[
\alpha,\ -\alpha,\ i\alpha,\ -i\alpha.
\]
因此分裂域为
\[
L=\Q(\alpha,i).
\]
先有 $[\Q(\alpha):\Q]=4$，而 $i\notin\Q(\alpha)$，所以
\[
[L:\Q]=8.
\]

Galois 群由
\[
\alpha\mapsto i\alpha,
\qquad i\mapsto i
\]
与复共轭
\[
i\mapsto -i
\]
生成，满足正方形对称群的关系，因此
\[
\Gal(L/\Q)\cong D_4.
\]
三个二次子域可由指数 $4$ 的正规子群给出，具体为
\[
\Q(i),
\qquad \Q(\sqrt2),
\qquad \Q(i\sqrt2).
\]
\end{solution}

\begin{problem}{Smith标准形的应用}
\hard
设整数矩阵
\[
A=\begin{pmatrix}2&4\\6&8\end{pmatrix}.
\]
求 $\Z^2/A\Z^2$ 的结构。
\end{problem}
\begin{solution}
对整数矩阵做 Smith 标准形。第一不变因子 $d_1$ 是所有条目的最大公因子：
\[
d_1=\gcd(2,4,6,8)=2.
\]
又
\[
\det A=16-24=-8,
\]
所以第二不变因子满足
\[
d_1d_2=|\det A|=8,
\qquad d_2=4.
\]
因此存在可逆整数矩阵 $U,V$ 使
\[
UAV=\operatorname{diag}(2,4).
\]
于是
\[
\Z^2/A\Z^2\cong \Z^2/\operatorname{diag}(2,4)\Z^2
\cong \Z/2\Z\oplus\Z/4\Z.
\]
\end{solution}

\begin{problem}{域的乘法群中的有限子群}
\hard
证明：域 $F$ 的乘法群 $F^\times$ 的任意有限子群都是循环群。
\end{problem}
\begin{solution}
设 $G\subset F^\times$ 是有限子群，记 $|G|=n$。则每个 $g\in G$ 都满足
\[
g^n=1,
\]
因为这是有限群中的 Lagrange 定理。于是 $G$ 的所有元素都是多项式
\[
x^n-1
\]
在域 $F$ 中的根。

设 $m$ 是 $G$ 中元素阶的最大值，取阶为 $m$ 的元素 $a$。则 $G$ 中任一元素的阶都整除 $m$，所以它们都是
\[
x^m-1
\]
的根。该多项式次数为 $m$，在域中至多有 $m$ 个根，因此
\[
|G|\le m.
\]
另一方面 $a$ 生成的循环子群已有 $m$ 个元素，所以 $m\le|G|$。故 $|G|=m$，于是 $G=\langle a\rangle$ 是循环群。
\end{solution}

\begin{problem}{Gauss引理}
\hard
证明：若 $f,g\in\Z[x]$ 都是 primitive 多项式，则乘积 $fg$ 仍是 primitive。并据此说明 primitive 多项式在 $\Z[x]$ 中不可约当且仅当它在 $\Q[x]$ 中不可约。
\end{problem}
\begin{solution}
设 $fg$ 的所有系数都被某素数 $p$ 整除。把 $f,g$ 模 $p$ 化到 $(\Z/p\Z)[x]$ 中，得到
\[
\bar f\,\bar g=0.
\]
但 $(\Z/p\Z)[x]$ 是整环，所以 $\bar f=0$ 或 $\bar g=0$，即 $f$ 或 $g$ 的所有系数都被 $p$ 整除。这与二者 primitive 矛盾。因此 $fg$ 仍 primitive。

若 primitive 多项式 $f\in\Z[x]$ 在 $\Q[x]$ 中可约，写成
\[
f=uv,
\qquad u,v\in\Q[x],
\]
清分母后可写成
\[
cf=f_1g_1,
\qquad f_1,g_1\in\Z[x].
\]
把内容提取出来并用上面的结论，可得到 $f$ 在 $\Z[x]$ 中也可约。反过来显然成立。所以 primitive 多项式在 $\Z[x]$ 与 $\Q[x]$ 中的不可约性等价。
\end{solution}

\begin{problem}{PID上有限生成模的结构定理}
\hard
陈述 PID 上有限生成模的结构定理，并用它分类所有形如
\[
M\cong \Z\oplus T
\]
且 $|T|=12$ 的有限生成 $\Z$-模。
\end{problem}
\begin{solution}
结构定理说：若 $R$ 是 PID，则任意有限生成 $R$-模都同构于
\[
R^r\oplus R/(d_1)\oplus\cdots\oplus R/(d_t),
\qquad d_1\mid d_2\mid\cdots\mid d_t,
\]
其中 $r\ge0$，这些不变因子唯一确定。

在这里 $R=\Z$，且自由部分秩为 $1$，扭部分 $T$ 是一个 $12$ 阶有限阿贝尔群。由前面的分类，$T$ 只有两种可能：
\[
\Z/12\Z,
\qquad
\Z/6\Z\oplus\Z/2\Z.
\]
因此所求模恰有两类：
\[
\Z\oplus\Z/12\Z,
\qquad
\Z\oplus\Z/6\Z\oplus\Z/2\Z.
\]
\end{solution}

\begin{problem}{中国剩余定理的一般形式}
\hard
设 $R$ 是交换环，$I_1,\dots,I_r$ 是两两互素的理想，即对 $i\neq j$ 有
\[
I_i+I_j=R.
\]
证明
\[
R/(I_1\cdots I_r)\cong R/I_1\times\cdots\times R/I_r.
\]
\end{problem}
\begin{solution}
考虑自然环同态
\[
\Phi:R\to R/I_1\times\cdots\times R/I_r,
\qquad
x\mapsto(x\bmod I_1,\dots,x\bmod I_r).
\]
其核为
\[
\bigcap_{i=1}^r I_i.
\]
对两两互素理想，可由归纳证明
\[
\bigcap_{i=1}^r I_i=I_1\cdots I_r.
\]
因此第一同构定理给出
\[
R/(I_1\cdots I_r)\cong \operatorname{im}\Phi.
\]

还需证明满射。先看 $r=2$：由 $I_1+I_2=R$，可取 $a\in I_1,b\in I_2$ 使
\[
a+b=1.
\]
于是给定 $(\bar x,\bar y)\in R/I_1\times R/I_2$，元素
\[
by+ax
\]
在模 $I_1$ 下同余于 $x$，模 $I_2$ 下同余于 $y$，故 $\Phi$ 满射。一般情形对 $r$ 归纳即可。

于是 $\operatorname{im}\Phi$ 就是整个直积，定理得证。
\end{solution}
